\documentclass[11pt]{amsart}

\usepackage[T1]{fontenc}
\usepackage{lmodern}
\usepackage[protrusion=true,expansion=false]{microtype}
\usepackage{amsmath,amssymb,amsthm,mathtools}
\usepackage[margin=2.5cm]{geometry}
\usepackage[colorlinks=true,linkcolor=blue,citecolor=blue,urlcolor=blue,
            hypertexnames=false]{hyperref}
\usepackage{setspace}

\newtheorem{mainresult}{Theorem}
\renewcommand{\themainresult}{\Alph{mainresult}}
\newtheorem{maincorollary}[mainresult]{Corollary}
\newtheorem{theorem}{Theorem}[section]
\newtheorem{proposition}[theorem]{Proposition}
\newtheorem{lemma}[theorem]{Lemma}
\newtheorem{corollary}[theorem]{Corollary}
\theoremstyle{remark}
\newtheorem{remark}[theorem]{Remark}


\newcommand{\R}{\mathbb R}
\newcommand{\C}{\mathbb C}
\newcommand{\T}{\mathbb T}
\newcommand{\HH}{\mathcal H}
\newcommand{\Sc}{\mathcal S}
\newcommand{\supp}{\operatorname{supp}}
\newcommand{\Ran}{\operatorname{Ran}}
\newcommand{\ac}{\mathrm{ac}}
\newcommand{\pp}{\mathrm{pp}}
\newcommand{\Id}{\mathrm{Id}}
\newcommand{\D}{\mathcal D}
\newcommand{\1}{\mathbf 1}
\newcommand{\loc}{\mathrm{loc}}
\setstretch{1.25}

\title{Phase retrieval for Schr\"odinger evolutions}
\author{Ben Pineau}
\author{João P. G. Ramos}
\author{Mitchell A. Taylor}
\date{}

\begin{document}
\raggedbottom

\begin{abstract}
We prove that two $L^2$ solutions of the one-dimensional Schr\"odinger equation with the same
real potential are equal up to one constant unimodular factor whenever their moduli agree.  The potential may depend on time and is
assumed locally to belong to
$L^1_tL^\infty_x+L^2_{t,x}$.

The result yields
phase retrieval for the cubic nonlinear Schr\"odinger equation in the global $L^2$ class.  The conclusion is moreover \emph{false} for the free equation in every spatial dimension at least two, even for Schwartz initial data.
\end{abstract}

\maketitle


\section{Introduction}

Let $\Omega\subset\R^n$ be a domain, and let $T>0$. We say that a function
$w:[0,T]\times\Omega\to\C$ solves a Schr\"odinger equation with potential
$V:[0,T]\times\Omega\to\R$ if it satisfies
\begin{equation}\label{eq:schr-pot}
i\partial_t w+\Delta w=Vw,
\end{equation}
in a suitable sense. Equation \eqref{eq:schr-pot} is of pivotal importance in quantum mechanics, where by adjusting the potential $V$ we may use it to describe the time evolution of the spatial distribution of particles. We refer the reader to Teschl~\cite{TeschlBook} for the theoretical background in physics of this equation.

Our question in this manuscript is tightly connected to quantum mechanics, and dates back to the work of Pauli~\cite{Pauli}. Indeed, let us consider first the problem of determining when it is possible to recover a (nice enough) function $f$ up to a unimodular constant from measurements of the form $|f|$, $|\widehat{f}|$. This is the so-called \emph{Pauli problem}. A negative answer was later given by Moroz and Perelomov~\cite{MorozPerelomov}, who showed that counterexamples are, as a matter of fact, abundant.

As it turns out, the Fourier transform is intimately connected to Schr\"odinger equations, and so is the Pauli problem: if we let $u$ be a solution to
\begin{equation}\label{eq:free-schr}
\begin{cases}
i\partial_t u(t,x)+\Delta u(t,x)=0 & \text{in }\R\times\R^n,\\
u(0,x)=f(x) & \text{on }\R^n.
\end{cases}
\end{equation}
Then, for $t\ne0$,
\[
u(t,x)=\frac{1}{(4\pi i t)^{n/2}}e^{i|x|^2/(4t)}
\widehat{\left(e^{i|\cdot|^2/(4t)}f\right)}
\left(\frac{x}{4\pi t}\right),
\]
where $\widehat h(\xi)=\int_{\R^n}e^{-2\pi i x\cdot\xi}h(x)\,dx$. In other words, if we are able to solve Pauli's problem for the corresponding chirped datum, we should be able to recover the solution $u$ to the free Schr\"odinger equation \eqref{eq:free-schr} above from its values $|u(0,\cdot)|$, $|u(t_0,\cdot)|$, and vice-versa. Regarding the problem in this particular setting, the counterexamples to Pauli's problem do show also here that recovery of the solution is not possible in general. This led the community to several follow-up questions, such as Wright's conjecture~\cite{JamingRathmair}, which states that there is a third unitary operator $T$ defined on $L^2$ such that any $f \in L^2(\R)$ may be recovered from $|f|,|Tf|, |\widehat{f}|$. Another closely related question~\cite{CarmeliHeinosaariSchultzToigo} is that of determining the minimal $N$ such that there are $N$ times $\{t_i\}_{i=1}^N$ for which any solution to \eqref{eq:free-schr} is determined up to multiplication by scalars from $\{|u(t_i,x)|:x\in\R,\ i=1,\dots,N\}$. To the best of our knowledge, both problems remain open.

As \eqref{eq:free-schr} is a particular case of \eqref{eq:schr-pot}, we may ask more generally if one may recover the solution $w$ to \eqref{eq:schr-pot} from the data $|w(0,\cdot)|$, $|w(t_0,\cdot)|$, where $t_0\in\R$ is a fixed time. In this direction, Jaming~\cite{Jaming2014} first showed that if one considers solutions to \eqref{eq:free-schr}, then phase retrieval is possible from $\{|u(t,x)|\}_{(t,x)\in\R\times\R}$ by connecting the problem to the fractional Fourier transform. Moreover, he later conjectured~\cite{Jaming2025} that the same should hold for \eqref{eq:schr-pot} under certain conditions on the potential $V$, and claimed that the same result should hold in higher dimensions.

The main purpose of this manuscript is to answer some of the latter questions on phase retrieval for solutions to \eqref{eq:schr-pot}. Our main result is the following:

\begin{mainresult}
\label{thm:time-potential}
Let $I_0\subset\R$ be a nonempty open interval. Suppose that the measurable
function $V:I_0\times\mathbb R\to\mathbb R$ has, on every bounded interval
$J\Subset I_0$, a decomposition
\[
 V=V_\infty+V_2,
 \qquad
 V_\infty\in L^1(J;L^\infty(\mathbb R)),
 \qquad
 V_2\in L^2(J\times\mathbb R).
\]
Let $u,v\in C(I_0;L^2(\mathbb R))$ be mild solutions of\eqref{eq:schr-pot}. If
\[
 |u(t,x)|=|v(t,x)|
 \quad\text{for almost every }(t,x)\in I_0\times\mathbb R,
\]
then $v=\zeta u$ on $I_0\times\mathbb R$ for some
$\zeta\in\mathbb T$.
\end{mainresult}

Here, a solution $w$ to \eqref{eq:schr-pot} is said to be ``mild'' if it satisfies the Duhamel formula. That is, for $t_0,t\in I_0$, we have
\[
 w(t)=e^{i(t-t_0)\partial_x^2}w(t_0)
 -i\int_{t_0}^t e^{i(t-s)\partial_x^2}(V(s)w(s))\,ds
\]
in the sense of distributions. The products of $w$ with each part of the potential satisfy
\[
 V_\infty w\in L^1(I;L^2),
 \qquad
 V_2w\in L^2(I;L^1)\subset L^{4/3}(I;L^1),
\]
by H\"older's inequality and the fact that $w\in C(I;L^2)$. The Duhamel integral above is then interpreted using this decomposition.

We remark that, in particular, the potentials in~\cite{Jaming2025} are all assumed to be time-\emph{independent}, and hence Theorem \ref{thm:time-potential} provides a positive answer to Question I.2 in~\cite{Jaming2025} for a wide class of potentials, namely any potential in $L^2(\R)+L^\infty(\R)$, which includes all $L^p(\R)$,
$2\leq p\leq\infty$.

The proof of Theorem \ref{thm:time-potential} starts with the endpoint Strichartz and local smoothing estimates. On each compact time interval,  we have
\[
 u,v\in L^4_tL^\infty_x,
 \qquad
 \chi u,\chi v\in L^2_tH^{1/2}_x
 \quad(\chi\in C_c^\infty(\R)).
\]
This is enough to justify the continuity equation
\[
 \partial_t|w|^2+2\partial_x\Im(\overline w\,\partial_xw)=0
 \]
in $L^2$. A regularization argument together with $|u| = |v|$ then gives
\[
 v\,\partial_xu-u\,\partial_xv=0
 \quad\text{in }\mathcal D'(\R_x)
\]
for almost every time. The next step is then to introduce the exterior product
\[
 F(t,x,y)=u(t,x)v(t,y)-v(t,x)u(t,y).
\]
Our goal is to show this vanishes identically, as this proves the asserted claim about linear dependence of the two solutions. In order to do so, note that it solves a two-dimensional Schr\"odinger equation with potential $V(t,x)+V(t,y)$. In the coordinates $r=(x+y)/\sqrt2$ and $s=(x-y)/\sqrt2$, the function $F$ is odd in $s$. Its Dirichlet trace at $s=0$ vanishes by oddness, and its normal trace vanishes by the Wronskian identity. The zero extension of $F$ from $s>0$ therefore solves the same equation without a boundary source.

Finally, a stripwise version of the Ionescu--Kenig Carleman estimate propagates this zero half-space through the whole $(r,s)$-plane. The $L^1_tL^\infty_x$ part of the potential is made small by shortening the time interval, while the restriction of the $L^2_{t,x}$ part to a strip of width $h$ has norm of order $h^{1/2}$. Taking $h$ small then allows us to iterate the estimate and gives $F=0$. Thus $u$ and $v$ are linearly dependent at one time, and equality of their moduli makes the proportionality constant unimodular, with uniqueness propagating the same phase throughout $I_0$.

If the potential and both solutions extend globally, uniqueness propagates the conclusion to every time.  The theorem also applies to nonlinear equations whenever equality of the moduli makes their nonlinear potentials
identical. This gives immediately the following consequence.

\begin{maincorollary}
\label{cor:cubic-NLS}
Let $\sigma\in\mathbb R$, let $I_0\subset\R$ be a nonempty open interval, and let
$u,v\in C(\mathbb R;L^2(\mathbb R))$ be global mild solutions of
\[
i\partial_t w+\partial_x^2w=\sigma|w|^2w.
\]
If $|u|=|v|$ almost everywhere on $I_0\times\mathbb R$, then
$v=\zeta u$ for every time and some $\zeta\in\mathbb T$.
\end{maincorollary}

As a final result, we prove that the prediction made by Jaming in~\cite{Jaming2025} regarding the extendability of the results for \eqref{eq:free-schr} to higher dimensions is \emph{false}:

\setcounter{mainresult}{5}
\begin{mainresult}
\label{thm:higher-dimensional-failure}
For every $d\ge2$ there are nonproportional functions
$f,g\in\mathcal S(\mathbb R^d)$ such that
\[
 |e^{it\Delta}f(x)|=|e^{it\Delta}g(x)|
 \quad\text{for every }(t,x)\in\mathbb R\times\mathbb R^d.
\]
\end{mainresult}

The counterexample is based on the Gaussian $\phi(x)=e^{-x^2/2}$ and its negative derivative $\psi(x)=xe^{-x^2/2}$. In dimension two, set
\[
 f=\phi\otimes\psi+i\psi\otimes\phi,
 \qquad
 g=\phi\otimes\psi-i\psi\otimes\phi.
\]
These are nonproportional Schwartz functions. If $A_t=e^{it\partial_x^2}\phi$, then $e^{it\partial_x^2}\psi=xA_t/(1+2it)$, and the tensor-product factorization of the free propagator gives
\begin{align*}
 e^{it\Delta}f(x_1,x_2)
 &=\frac{A_t(x_1)A_t(x_2)}{1+2it}(x_2+ix_1),\\
 e^{it\Delta}g(x_1,x_2)
 &=\frac{A_t(x_1)A_t(x_2)}{1+2it}(x_2-ix_1).
\end{align*}
The two expressions have the same modulus for every point and every time. For $d>2$, tensor both functions with a Gaussian in the remaining variables.

This article is organized as follows. In Section 2, we prove Theorem \ref{thm:time-potential}. In Section 3, we deduce Corollary \ref{cor:cubic-NLS} as a consequence of the preceding result. Section 4 is then dedicated to the counterexample construction of Theorem \ref{thm:higher-dimensional-failure}. Finally, we gather, in Section 5 below, some other results obtained in parallel to the ones here. These are results which are mostly superseded/covered by the ones in the present paper, but since their techniques are interesting and beautiful on their own, we decided to keep them in a separate repository for independent consultation by the interested reader.

\subsection*{LLM Usage.} Large language models have been of great importance in the production of this manuscript. Below is a brief narration of how they contributed to the current work. 

We started from the first idea of considering the continuity equation $ \partial_t|w|^2+2\partial_x\Im(\overline w\,\partial_xw)=0$, together with integration by parts, in order to show that, if we had $w(t,x) = |w(t,x)| \cdot e^{i \phi(t,x)},$ with $\phi$ smooth enough on each connected component where $w$ does not vanish, then $\phi$ must be constant on each such component. Our initial argument then ran through classical works (see e.g. \cite{Masuda}) to obtain uniqueness of the phase overall. 

This argument, however, only works in high enough regularity. In order to remove such regularity, we needed to get rid of the assumptions on the phase. At that point, we started exploring with the aid of GPT 5.4/Claude Opus 4.6. It gave us several different directions, many of which are described in the end of this manuscript, such as ones based on scattering theory. 

By exploiting and exhausting several paths towards the general $L^2$ case with GPT 5.4 and 5.5, as well as Claude Opus 4.7 and 4.8, we had the idea of working at the level of currents. By doing so, we were naturally led to the idea of lifting our equation one dimension higher. The details on how to finish from there, such as applying specific unique continuation lemmas, were then provided by interacting with several different LLMs. Finally, Claude Fable 5 was also used in order to obtain a Lean verification of Corollary \ref{cor:cubic-NLS}; see \href{https://github.com/joaopgramos95/Schrodinger-PR/tree/main/SchrodingerPR_verified}{this link}. 

Most of the other attempts and results through different methods have been written down and are kept \href{https://github.com/joaopgramos95/Schrodinger-PR/tree/main/Article}{here}. Apart from that note -- which is mostly LLM-written --, the results and proofs here have all been written and verified by humans. We take full responsibility for the contents of both manuscripts. 

\section*{Acknowledgements}

The authors are thankful to Philippe Jaming and Jaume de Dios for several conversations about the topic during the elaboration of the manuscript. 

J. P. G. R.  was supported by the FCT through project SHADE
(2023.17881.ICDT), DOI \allowbreak \texttt{10.54499/}
\texttt{2023.17881.ICDT}. He was also supported  by FAPERJ through  JCNE grant
no.~SEI-260003/ \allowbreak 020475/2025, and by Instituto Serrapilheira through
grant Serra-R-2510-59765.

\section{Proof of Theorem \ref{thm:time-potential}}

We start with the following local form of the half-space propagation argument of
Ionescu--Kenig~\cite{IonescuKenig}.

\begin{proposition}
\label{prop:strip-propagation}
Let $J$ be a compact interval, $z=(z_1,z_2)\in\R^2$, and
$\mathcal H=i\partial_t+\Delta_z$.  Put
\begin{align*}
 X&=L^1(J;L^2(\R^2))+L^{4/3}(J\times\R^2),\\
 X'&=L^\infty(J;L^2(\R^2))\cap L^4(J\times\R^2),\\
 Y&=L^1(J;L^\infty(\R^2))+L^2(J\times\R^2).
\end{align*}
The sum norm on $X$ and $Y$ is the infimum over all decompositions into the
two displayed summands; the intersection norm on $X'$ is the maximum of the
two norms.  There is an absolute $\varepsilon_*>0$ with the following
property.  Suppose
$Z\in C(J;L^2)\cap X'$, $\mathcal HZ=WZ\in X$, and
$Z=0$ on $J\times\{z_1>b\}$.  If for some $h>0$
\[
 \sup_{a\le b-h}
 \|\1_{\{a<z_1<a+h\}}W\|_Y\le\varepsilon_*,
 \tag{1.3}\label{eq:strip-smallness}
\]
then $Z=0$ on $J\times\R^2$.
\end{proposition}

\begin{proof}
The multiplier estimate used in the two-dimensional Schr\"odinger Carleman
inequality follows from two applications of H\"older's inequality.  If
$A=A_\infty+A_2$ with
$A_\infty\in L^1_tL^\infty_z$ and $A_2\in L^2_{t,z}$, then
\begin{align*}
 \|A_\infty B\|_{L^1_tL^2_z}
 &\le\|A_\infty\|_{L^1_tL^\infty_z}
       \|B\|_{L^\infty_tL^2_z},\\
 \|A_2B\|_{L^{4/3}_{t,z}}
 &\le\|A_2\|_{L^2_{t,z}}\|B\|_{L^4_{t,z}}.
\end{align*}
Taking the infimum over decompositions of $A$ gives
\[
 \|AB\|_X\le
 \|A\|_Y\|B\|_{X'}.
 \tag{1.4}\label{eq:XY-multiplier}
\]
Apply the bounded-weight estimate in
\cite[Lemma~9.1 and the proof of Theorem~2.5]{IonescuKenig} after translating
the plane $z_1=b$ to $z_1=0$.  Its potential term is bounded by
\eqref{eq:XY-multiplier} and is absorbed when its $Y$ norm is at most
$\varepsilon_*$.  After the potential term has been absorbed, the Carleman parameter may tend
to infinity.  The weighted estimate then shows that a solution which vanishes
on $z_1>b$ also vanishes on $z_1>b-h$.

For precision, the time-dependent translation used in the cited proof is
$z_1=y_1+w_\delta(t)$, where $w_\delta=b$ near the endpoints of $J$ and
$w_\delta=b-h$ away from them.  Conjugation by
$e^{-iw_\delta'(t)y_1/2}$ removes the first-order term.  The remaining affine
term $w_\delta''(t)y_1/2$ is controlled by the coercive bounded weight in that
lemma.  No potential outside the strip $b-h<z_1<b$ occurs because $Z$ already
vanishes to its right.  Letting the Carleman parameter tend to infinity and
then $\delta\downarrow0$ gives the asserted enlargement of the zero set.
Iteration over the planes $z_1=b-jh$, $j\ge1$, gives
\[
 Z=0\quad\text{on }\{z_1>b-jh\}\times J
 \quad(j=0,1,2,\ldots).
\]
The union of these half-spaces is $\R^2$, which proves $Z=0$.  
\end{proof}

We next recover the normal Cauchy data needed to create a half-space solution.

\begin{lemma}
\label{lem:endpoint-wronskian}
Under the hypotheses of Theorem~\ref{thm:time-potential}, on every bounded
interval $I\Subset I_0$ one has
\[
 u,v\in L^4(I;L^\infty(\R)),
 \qquad
 \chi u,\chi v\in L^2(I;H^{1/2}(\R))
 \quad(\chi\in C_c^\infty(\R)).
 \tag{1.5}\label{eq:linear-endpoint-regularity}
\]
Moreover, for almost every $t\in I_0$,
\[
 \overline u\,\partial_xu=\overline v\,\partial_xv,
 \qquad
 v\,\partial_xu-u\,\partial_xv=0
 \quad\text{in }\D'(\R_x),
 \tag{1.6}\label{eq:current-wronskian}
\]
where the first product is interpreted by
$H^{-1/2}$--$H^{1/2}$ duality.
\end{lemma}

\begin{proof}
Fix $I=[a,b]$.  The one-dimensional admissible endpoint is
$(q,r)=(4,\infty)$, since $2/q+1/r=1/2$.  The homogeneous and inhomogeneous
Strichartz estimates~\cite{KeelTao} give
\begin{align*}
 \|e^{i(t-a)\partial_x^2}w(a)\|_{L^4_tL^\infty_x(I)}
 &\lesssim\|w(a)\|_2,\\
 \left\|\int_a^t e^{i(t-s)\partial_x^2}F(s)\,ds
 \right\|_{L^4_tL^\infty_x(I)}
 &\lesssim\|F\|_{L^1_tL^2_x+L^{4/3}_tL^1_x}.
\end{align*}
Apply these estimates to the Duhamel formula.  For either solution $w \in \{u,v\}$, this
gives
\begin{align*}
 \|w\|_{L^4_tL^\infty_x(I)}
 \lesssim_I{}&\|w\|_{L^\infty_tL^2_x(I)}
 +\|V_\infty\|_{L^1_tL^\infty_x(I)}
       \|w\|_{L^\infty_tL^2_x(I)}\\
 &+|I|^{1/4}\|V_2\|_{L^2_{t,x}(I\times\R)}
       \|w\|_{L^\infty_tL^2_x(I)}.
\end{align*}
For the last term, H\"older's inequality in $x$ gives
\[
 \|V_2(t)w(t)\|_1\le\|V_2(t)\|_2\|w(t)\|_2.
\]
The function $t\mapsto\|V_2(t)\|_2$ belongs to $L^2(I)$, and
$L^2(I)\subset L^{4/3}(I)$ with norm at most $|I|^{1/4}$.  This proves the
displayed Strichartz bound.  Since $L^4(I)\subset L^2(I)$ on a bounded
interval, we also obtain
\begin{align*}
 \|V_2w\|_{L^1_tL^2_x}
 &\le \|V_2\|_{L^2_{t,x}}\|w\|_{L^2_tL^\infty_x}<\infty,\\
 \|V_\infty w\|_{L^1_tL^2_x}
 &\le \|V_\infty\|_{L^1_tL^\infty_x}
       \|w\|_{L^\infty_tL^2_x}<\infty.
\end{align*}
The local-smoothing estimate~\cite{VegaVisciglia} states that, for
$\chi\in C_c^\infty(\R)$,
\[
 \|\chi e^{it\partial_x^2}g\|_{L^2(I;H^{1/2})}
 \le C_{I,\chi}\|g\|_2.
\]
Apply this estimate to the homogeneous term in the Duhamel formula.  For the
inhomogeneous term, apply it for each fixed forcing time $s$ and then use
Minkowski's integral inequality.  The two preceding $L^1_tL^2_x$ bounds give
\[
 \left\|\chi\int_a^t e^{i(t-s)\partial_x^2}(Vw)(s)\,ds
 \right\|_{L^2_tH^{1/2}_x}
 \lesssim_{I,\chi}\|Vw\|_{L^1_tL^2_x}.
\]
This proves \eqref{eq:linear-endpoint-regularity}.  For the remainder of the
proof, take a bounded interval $I$ with $\overline I\subset I_0$.

For $w\in H^{1/2}_{\loc}\cap L^\infty_{\loc}$ define
\[
 \langle\overline w\,\partial_xw,\phi\rangle
 :=\langle\partial_xw,\phi\overline w\rangle_{H^{-1/2},H^{1/2}},
 \qquad \phi\in C_c^\infty(\R).
\]
This pairing is well-defined. Indeed, after inserting a cutoff equal to one on
$\supp\phi$, multiplication by $\phi$ maps
$H^{1/2}\cap L^\infty$ into $H^{1/2}$, while differentiation maps
$H^{1/2}$ into $H^{-1/2}$.  Its time integral is finite because the local
$H^{1/2}$ norm belongs to $L^2(I)$.

With that object being well-defined, the continuity equation
\[
 \partial_t|w|^2+2\partial_x\Im(\overline w\,\partial_xw)=0
 \tag{1.7}\label{eq:continuity-general}
\]
holds in distributions. Let now $\phi\in C_c^\infty(\R)$ and
$a(x)=\int_{-\infty}^x\phi(y)\,dy$. It follows that $a$ is bounded and $a'=\phi$. Although $a$ need not have compact support, it is a bounded Lipschitzmultiplier and $a'=\phi$ is compactly supported.  One may therefore test the
weak equation \eqref{eq:continuity-general} directly against a regularization of
$\eta(t)a(x)\overline{w(t,x)}$.  Multiplication by $a$ is bounded on
$H^{1/2}$, the potential contribution is real and cancels after taking the
imaginary part, and all remaining pairings are integrable in time by
\eqref{eq:linear-endpoint-regularity}.  Approximation by smooth functions
then gives
\[
 -\int_I\eta'(t)\int_\R a(x)|w(t,x)|^2\,dx\,dt
 =2\int_I\eta(t)
 \langle\Im(\overline w\,\partial_xw),\phi\rangle\,dt.
\]
Equivalently,
\[
 \frac d{dt}\int_\R a|w|^2
 =2\langle\Im(\overline w\,\partial_xw),\phi\rangle
 \tag{1.8}\label{eq:localized-continuity-general}
\]
in $\D'(I)$.  The left-hand sides for $u$ and $v$ agree, since their absolute values agree.  Hence
\[
 \int_I\eta(t)
 \langle\Im(\overline u\,\partial_xu-\overline v\,\partial_xv),\phi\rangle\,dt=0
\]
for every $\eta\in C_c^\infty(I)$ and every spatial test function $\phi$.
Choose a countable dense family of spatial test functions.  After removing
one null set of times, the current identity holds for every member of this
family and therefore, by continuity of the distributional pairing, for every
$\phi$.  Their real parts agree because the distributional product rule gives
\[
 2\Re(\overline w\,\partial_xw)=\partial_x|w|^2.
\]
This rule follows by approximating $w$ in local $H^{1/2}$ while retaining its
local $L^\infty$ bound, or directly from the Gagliardo product estimate.  We
have proved the first identity in \eqref{eq:current-wronskian} for almost every
time.

We now prove the second identity.  The
following one-dimensional fact is used at almost every time.  Suppose
$p,q\in H^{1/2}_{\loc}\cap L^\infty_{\loc}$,
$|p|=|q|$, and
$\overline p\,\partial_xp=\overline q\,\partial_xq$ as quadratic-form
distributions.  After multiplying both functions by a common smooth cutoff which equals one
near the test-function support, assume
$p,q\in H^{1/2}\cap L^\infty$.  The hypotheses are preserved: the additional
term in each logarithmic derivative is the same because $|p|=|q|$.  Put
\[
 \rho=|p|^2=|q|^2,
 \quad s_\varepsilon=\frac{\rho}{\rho+\varepsilon},
 \quad h_\varepsilon=\frac{pq}{\rho+\varepsilon}.
\]
The Gagliardo seminorm is
\[
 [a]_{H^{1/2}}^2=\iint_{\R^2}
 \frac{|a(x)-a(y)|^2}{|x-y|^2}\,dx\,dy.
\]
It shows directly that a Lipschitz map fixing zero acts continuously on
$H^{1/2}$, and that $H^{1/2}\cap L^\infty$ is an algebra.  Hence
$h_\varepsilon$, $s_\varepsilon p$, and $s_\varepsilon q$ are admissible
localized $H^{1/2}$ multipliers.  Since
\[
 \overline p h_\varepsilon=s_\varepsilon q,
 \qquad
 \overline q h_\varepsilon=s_\varepsilon p,
\]
testing the logarithmic-derivative identity against
$h_\varepsilon\phi$ gives
\[
 \langle\partial_xp,s_\varepsilon q\phi\rangle
 -\langle\partial_xq,s_\varepsilon p\phi\rangle=0.
\]
Thus, letting $T_\varepsilon(z)=\varepsilon z/(|z|^2+\varepsilon)$,
\[
 \langle q\,\partial_xp-p\,\partial_xq,\phi\rangle
 =\langle\partial_xp,T_\varepsilon(q)\phi\rangle
  -\langle\partial_xq,T_\varepsilon(p)\phi\rangle.
\]
The maps $T_\varepsilon$ are uniformly Lipschitz, vanish at zero, and satisfy
$|T_\varepsilon(z)|\le\sqrt\varepsilon/2$.  Thus they converge pointwise to
zero.  Moreover,
\[
 \frac{|T_\varepsilon(p(x))-T_\varepsilon(p(y))|^2}{|x-y|^2}
 \le C\frac{|p(x)-p(y)|^2}{|x-y|^2},
\]
with $C$ independent of $\varepsilon$.  The Gagliardo formula and dominated
convergence give
$T_\varepsilon(p),T_\varepsilon(q)\to0$ in $H^{1/2}_{\loc}$.  After
multiplication by the fixed test function $\phi$, duality gives
\begin{align*}
 |\langle\partial_xp,T_\varepsilon(q)\phi\rangle|
 &\le \|\partial_xp\|_{H^{-1/2}(K)}
       \|T_\varepsilon(q)\phi\|_{H^{1/2}},\\
 |\langle\partial_xq,T_\varepsilon(p)\phi\rangle|
 &\le \|\partial_xq\|_{H^{-1/2}(K)}
       \|T_\varepsilon(p)\phi\|_{H^{1/2}},
\end{align*}
where $K$ contains $\supp\phi$.  Both right-hand sides above tend to zero.  It follows thus by passing to
the limit that $q\,\partial_xp-p\,\partial_xq=0$. The first identity in
\eqref{eq:current-wronskian} then follows at once. 
\end{proof}

We are now ready to prove our main result. 

\begin{proof}[Proof of Theorem~\ref{thm:time-potential}]
Fix a bounded interval $I$ with $\overline I\subset I_0$.  On this interval
set $\rho=|u|^2=|v|^2$ and define the exterior product
\[
 F(t,x,y)=u(t,x)v(t,y)-v(t,x)u(t,y).
 \tag{1.9}\label{eq:exterior-product}
\]
The two tensor products in \eqref{eq:exterior-product} belong to
$C(I;L^2(\R^2))$, as
$\|a\otimes b\|_{L^2(\R^2)}=\|a\|_2\|b\|_2$.  Apply
$i\partial_t+\partial_x^2+\partial_y^2$ to each term.  For instance,
\begin{align*}
 &(i\partial_t+\partial_x^2+\partial_y^2)
   (u(t,x)v(t,y))\\
 &\qquad=V(t,x)u(t,x)v(t,y)+V(t,y)u(t,x)v(t,y).
\end{align*}
The identical calculation for $v(t,x)u(t,y)$ gives
\[
 i\partial_tF+\partial_x^2F+\partial_y^2F=(V(t,x)+V(t,y))F
 \tag{1.10}\label{eq:exterior-equation}
\]
in distributions.  The calculation is justified by inserting smooth
approximations in the mild equation and passing to the limit in the spaces
established below.  On each bounded interval $I$,
\[
 F\in C(I;L^2(\R^2))\cap L^4(I\times\R^2).
 \tag{1.11}\label{eq:exterior-class}
\]
Indeed, interpolation between $L^\infty_tL^2_x$ and
$L^4_tL^\infty_x$ with equal weights gives
\[
 \|w\|_{L^8_tL^4_x}
 \le\|w\|_{L^\infty_tL^2_x}^{1/2}
     \|w\|_{L^4_tL^\infty_x}^{1/2}.
\]
Consequently
\begin{align*}
 \|u(t,x)v(t,y)\|_{L^4_{t,x,y}}^4
 &=\int_I\int_\R\int_\R|u(t,x)|^4|v(t,y)|^4\,dx\,dy\,dt\\
 &=\int_I\|u(t)\|_4^4\|v(t)\|_4^4\,dt\\
 &\le\|u\|_{L^8_tL^4_x}^4\|v\|_{L^8_tL^4_x}^4<\infty.
\end{align*}
The same estimate holds for $v(t,x)u(t,y)$.
Moreover, the right side of \eqref{eq:exterior-equation} belongs to
$L^1(I;L^2(\R^2))$.  For example, Fubini's theorem gives
\begin{align*}
 &\|V_2(t,x)u(t,x)v(t,y)\|_{L^2_{x,y}}\\
 &\qquad=\|V_2(t)u(t)\|_{L^2_x}\|v(t)\|_{L^2_y}
 \le\|V_2(t)\|_2\|u(t)\|_\infty\|v(t)\|_2.
\end{align*}
The last expression is integrable in $t$ by Cauchy--Schwarz, because
$V_2\in L^2_tL^2_x$, $u\in L^2_tL^\infty_x$, and
$v\in L^\infty_tL^2_x$.  For $V_\infty$, use
\[
 \|V_\infty(t,x)u(t,x)v(t,y)\|_2
 \le\|V_\infty(t)\|_\infty\|u(t)\|_2\|v(t)\|_2
\]
and $V_\infty\in L^1_tL^\infty_x$.  The terms with $x$ and $y$ interchanged follow by identical means.

Let us now introduce orthogonal coordinates:
\[
 r=\frac{x+y}{\sqrt2},
 \qquad s=\frac{x-y}{\sqrt2}.
\]
Then
\[
 i\partial_tF+\partial_r^2F+\partial_s^2F=QF,
 \qquad
 Q(t,r,s)=V\!\left(t,\frac{r+s}{\sqrt2}\right)
          +V\!\left(t,\frac{r-s}{\sqrt2}\right).
 \tag{1.12}\label{eq:rotated-exterior}
\]
Interchanging $x$ and $y$ changes the sign of $F$.  In the new coordinates
this interchange fixes $r$ and sends $s$ to $-s$.  Therefore
\[
 F(t,r,-s)=-F(t,r,s).
\]
We claim that both of its Cauchy traces at $s=0$ vanish.  For $\psi\in C_c^\infty(I\times\R_r)$ set
\[
 H_\psi(s)=\langle F(\cdot,\cdot,s),\psi\rangle_{t,r}.
\]
Cauchy--Schwarz in $(t,r)$ and Fubini's theorem give
$H_\psi\in L^2_{\loc}(\R_s)$.  Pairing
\eqref{eq:rotated-exterior} with $\psi$ and integrating by parts in $t$ and
$r$ gives, in distributions of $s$,
\[
 H_\psi''(s)
 =\langle QF,\psi\rangle_{t,r}
 +\langle F,i\partial_t\psi-\partial_r^2\psi\rangle_{t,r}.
 \tag{1.12a}\label{eq:scalar-trace-equation}
\]
For a compact $s$-interval, Minkowski's inequality bounds the first term in
$L^2_s$ by a constant times $\|QF\|_{L^1_tL^2_{r,s}}$; the second is bounded
by $\|F\|_{L^1_tL^2_{r,s}}$.  Thus $H_\psi''\in L^2_{\loc}(\R_s)$.  The
one-dimensional Fourier characterization of Sobolev spaces now gives
$H_\psi\in H^2_{\loc}(\R_s)$.  Such a function and its first derivative have
continuous representatives, so both traces at $s=0$ exist.  Oddness gives
$H_\psi(0)=0$.

For $z=(r-s)/\sqrt2$,
\begin{align*}
 H_\psi(s)=\sqrt2\iint
 &[u(t,z+\sqrt2s)v(t,z)-v(t,z+\sqrt2s)u(t,z)]\\
 &\hspace{28mm}\times\psi(t,\sqrt2z+s)\,dz\,dt.
\end{align*}
For $p\in H^{1/2}$,
\[
 \frac{p(\cdot+h)-p}{h}\longrightarrow \partial_xp
 \quad\text{in }H^{-1/2}.
\]
Indeed, on the Fourier-transform side the multiplier is
$(e^{ih\xi}-1)/h$, which converges pointwise to $i\xi$ and is bounded by
$|\xi|$.  Dominated convergence with the $H^{1/2}$ weight proves the claim.
After inserting cutoffs supported near the projection of $\supp\psi$, pair
this convergence with the other local $H^{1/2}$ factor.  The difference
quotients are uniformly bounded from $H^{1/2}$ to $H^{-1/2}$, while the two
local $H^{1/2}$ norms belong to $L^2_t$.  Cauchy--Schwarz in time therefore
justifies passage to the limit under the temporal integral.  The term obtained by
differentiating $\psi(t,\sqrt2z+s)$ vanishes at $s=0$ because the bracket in
the preceding display vanishes there.  Lemma~\ref{lem:endpoint-wronskian}
then gives
\[
 H_\psi'(0)=2\iint
 \bigl(v\,\partial_xu-u\,\partial_xv\bigr)(t,z)
 \psi(t,\sqrt2z)\,dz\,dt=0.
\]
Thus the zero extension
\[
 G(t,r,s)=\1_{\{s>0\}}F(t,r,s)
\]
solves the same equation on all of $I\times\R^2$.  Indeed, the distributional
jump formula is
\[
 \partial_s^2(\1_{\{s>0\}}F)
 =\1_{\{s>0\}}\partial_s^2F
  +\delta_{s=0}(\partial_sF)|_{s=0}+\delta'_{s=0}F|_{s=0}.
\]
Both boundary distributions vanish by the two trace identities proved above.  Therefore, it follows that $G$ satisfies
\[
 i\partial_tG+\partial_r^2G+\partial_s^2G=QG.
\]
Multiplication by a half-space characteristic function is a bounded operator
on $L^2$, so $G\in C(I;L^2)$.  It also preserves the $L^4$ bound, and
$QG\in L^1_tL^2$ by the estimate already proved for $QF$. It remains hence to apply Proposition~\ref{prop:strip-propagation}.  Choose a compact
subinterval $J\subset I$ with sufficiently small length, so that
\[
 2\|V_\infty\|_{L^1(J;L^\infty)}<\varepsilon_*/2.
\]
For $S_{a,h}=\{(r,s):a<s<a+h\}$, Fubini's theorem and the substitution
$x=(r+s)/\sqrt2$, for which $dr=\sqrt2\,dx$, give
\begin{align*}
 &\left\|\1_{S_{a,h}}
 V_2\!\left(t,\frac{r+s}{\sqrt2}\right)\right\|_{L^2_{t,r,s}}^2\\
 &\quad=\int_J\int_a^{a+h}\int_\R
 \left|V_2\!\left(t,\frac{r+s}{\sqrt2}\right)\right|^2dr\,ds\,dt\\
 &\quad=\sqrt2h\int_J\int_\R|V_2(t,x)|^2\,dx\,dt
 =\sqrt2h\|V_2\|_{L^2(J\times\R)}^2.
 \tag{1.13}\label{eq:ridge-strip}
\end{align*}
and the same identity holds for the second ridge.  Choose $h>0$ so that their
combined $L^2$ norm is below $\varepsilon_*/2$.  After reflecting $s$ to make
the zero half-space $z_1>0$, \eqref{eq:ridge-strip} and
the hypotheses on the potential show that \eqref{eq:strip-smallness} holds.
Proposition~\ref{prop:strip-propagation} then gives at once that $G=0$.  By oddness, it follows then that $F=0$ almost everywhere on $J\times\R^2$.  Finally, since 
$F\in C(J;L^2(\R^2))$, it follows that $F(t)=0$ in $L^2(\R^2)$ for every
$t\in J$.

Choose now $t_0\in J$.  If $u(t_0)=0$, since $|u(t_0)| = |v(t_0)|$, it follows at once that $v(t_0)=0$. We then recall that if two mild solutions have the same value at
one endpoint of a sufficiently short interval $K$, their difference $w$
satisfies the homogeneous Duhamel equation.  The homogeneous and inhomogeneous
Strichartz estimates, together with the estimates used in
Lemma~\ref{lem:endpoint-wronskian}, give
\[
 \|w\|_{L^\infty_tL^2_x(K)\cap L^4_tL^\infty_x(K)}
 \le C\bigl(
 \|V_\infty\|_{L^1_tL^\infty_x(K)}
 +|K|^{1/4}\|V_2\|_{L^2_{t,x}(K)}\bigr)
 \|w\|_{L^\infty_tL^2_x(K)\cap L^4_tL^\infty_x(K)}.
\]
By absolute continuity of the two potential norms, $K$ may be chosen so that the coefficient on the right is less than one.  Hence $w=0$ on $K$.
Partitioning each compact subinterval of $I_0$ into finitely many such pieces
extends uniqueness forward and backward from $t_0$.  Since multiplication by a
constant commutes with the linear equation, $v$ and $\zeta u$ have the same
initial value at $t_0$ and therefore agree throughout $I_0$.

Suppose instead that $u(t_0)\ne0$.  Choose $x_0$ in a full-measure
set on which $u(t_0,x_0)\ne0$ and the identity
$F(t_0,x_0,y)=0$ holds for almost every $y$.  Then
\[
 v(t_0,y)=\frac{v(t_0,x_0)}{u(t_0,x_0)}u(t_0,y)
 \quad\text{for almost every }y.
\]
Thus $v(t_0)=\zeta u(t_0)$ for a scalar $\zeta$.  The hypotheses then directly imply that $|\zeta|=1$.
\end{proof}

\section{Cubic NLS: Proof of Corollary \ref{cor:cubic-NLS}}

We prove Corollary~\ref{cor:cubic-NLS}.  Its equation is
\[
 i\partial_t w+\partial_x^2w=\sigma|w|^2w.
 \tag{2.1}\label{eq:cubic-NLS}
\]
The mild solutions are the standard $L^2$ Strichartz solutions obtained from
the Duhamel formula.

\begin{proof}[Proof of Corollary~\ref{cor:cubic-NLS}]
Fix a bounded interval $I$.  A pair $(q,r)$ is Schr\"odinger-admissible in
one dimension if
\[
 \frac2q+\frac1r=\frac12,
 \qquad 2\le q,r\le\infty.
\]
The pair $(8,4)$ is admissible.  On a time interval $K$, the homogeneous and
inhomogeneous Strichartz estimates~\cite{KeelTao} give
\[
 \left\|e^{i(t-t_0)\partial_x^2}w(t_0)
 -i\sigma\int_{t_0}^te^{i(t-s)\partial_x^2}|w(s)|^2w(s)\,ds
 \right\|_{L^\infty_tL^2_x\cap L^8_tL^4_x}
 \lesssim \|w(t_0)\|_2
 +|\sigma|\,\||w|^2w\|_{L^{8/7}_tL^{4/3}_x}.
\]
For the nonlinear term, H\"older's inequality in space and the finite-measure
embedding in time give
\begin{align*}
 \||w|^2w\|_{L^{8/7}_tL^{4/3}_x(K)}
 &\le |K|^{1/2}
       \||w|^2w\|_{L^{8/3}_tL^{4/3}_x(K)}\\
 &=|K|^{1/2}\|w\|_{L^8_tL^4_x(K)}^3.
 \tag{2.2}\label{eq:NLS-fixed-point-estimate}
\end{align*}
The pointwise inequality
\[
 \bigl||z_1|^2z_1-|z_2|^2z_2\bigr|
 \le C(|z_1|^2+|z_2|^2)|z_1-z_2|
\]
and the same H\"older estimate give
\begin{align*}
 &\||w_1|^2w_1-|w_2|^2w_2\|_{L^{8/7}_tL^{4/3}_x(K)}\\
 &\qquad\le C|K|^{1/2}
 \bigl(\|w_1\|_{L^8_tL^4_x(K)}^2+\|w_2\|_{L^8_tL^4_x(K)}^2\bigr)
 \|w_1-w_2\|_{L^8_tL^4_x(K)}.
\end{align*}
Thus the Duhamel map is a contraction on a ball in
$C(K;L^2)\cap L^8(K;L^4)$ when $K$ is sufficiently short.  This is the local
$L^2$ construction referred to in the theorem.  For smooth solutions,
multiplication by $\overline w$, integration in $x$, and taking imaginary
parts give
\[
 \frac d{dt}\|w(t)\|_2^2=0.
\]
Approximation of the initial data and the $C_tL^2_x$ convergence in the fixed
point construction extend mass conservation to $L^2$ solutions.  Since the
length of the contraction interval depends only on $|\sigma|\|w(t_0)\|_2^2$,
the local solution can be iterated for all time.  In particular, covering the
bounded interval $I$ by finitely many local intervals gives
\[
 u,v\in L^8(I;L^4(\R)).
 \tag{2.3}\label{eq:NLS-L8L4}
\]
The inhomogeneous term belongs to $L^{8/7}_tL^{4/3}_x$ by
\eqref{eq:NLS-fixed-point-estimate}.  The pair $(8,4)$ is admissible, so its dual
pair is $(8/7,4/3)$.  The inhomogeneous Strichartz estimate from this dual
pair to the admissible endpoint $(4,\infty)$ gives
\[
 u,v\in L^4(I;L^\infty(\R)).
 \tag{2.4}\label{eq:NLS-endpoint}
\]
Finally, \eqref{eq:NLS-L8L4} and the finite length of $I$ imply
\[
 \|u\|_{L^4(I\times\R)}
 \le |I|^{1/8}\|u\|_{L^8_tL^4_x},
\]
and the same estimate holds for $v$.

On $I_0\times\R$, set $\rho=|u|^2=|v|^2$ and
$V=\sigma\rho$.  For every bounded $J\Subset I_0$, the preceding estimate
gives
\[
 \|V\|_{L^2(J\times\R)}
 =|\sigma|\|u\|_{L^4(J\times\R)}^2<\infty.
\]
Because $\sigma$ and $\rho$ are real, $V$ is real.  On $I_0$ the two
nonlinear equations are therefore the same linear equation:
\[
 i\partial_tu+\partial_x^2u=Vu,
 \qquad
 i\partial_tv+\partial_x^2v=Vv.
\]
Theorem~\ref{thm:time-potential} gives $v=\zeta u$ on $I_0$ for some
$\zeta\in\T$.  Choose $t_0\in I_0$.  The fixed-point difference estimate
above gives uniqueness for the cubic equation in
$C_tL^2_x\cap L^8_tL^4_x$.  Since the two nonlinear solutions agree up to
$\zeta$ at $t_0$, gauge covariance and uniqueness give
$v=\zeta u$ for every time.
\end{proof}

\section{Failure in dimensions at least two: Proof of Theorem \ref{thm:higher-dimensional-failure}}

As a final contribution, we build a simple family of counterexamples showing Theorem~\ref{thm:higher-dimensional-failure}. This, incidentally, provides both a \emph{negative} answer to the higher-dimensional assertion in \cite{Jaming2025}, and an example showing that the one-dimensional case of the four-time problem mentioned previously is genuinely special, as the same result in higher dimensions fails even for a continuum of times. 

\begin{proof}[Proof of Theorem~\ref{thm:higher-dimensional-failure}]
It suffices to work in $d=2$.  Let
\[
 \phi(x)=e^{-x^2/2},
 \qquad \psi(x)=-\phi'(x)=xe^{-x^2/2},
\]
and define
\begin{align*}
 a(x_1,x_2)&=\phi(x_1)\psi(x_2),\\
 b(x_1,x_2)&=\psi(x_1)\phi(x_2),\\
 f&=a+ib,
 \qquad g=a-ib.
\end{align*}
Since
\[
 f(x_1,x_2)=e^{-(x_1^2+x_2^2)/2}(x_2+ix_1),
 \qquad
 g(x_1,x_2)=e^{-(x_1^2+x_2^2)/2}(x_2-ix_1),
\]
both functions are Schwartz.  They are not scalar multiples.  If $f=cg$,
then setting $x_1=0$ and $x_2\ne0$ gives $c=1$, whereas setting $x_2=0$ and
$x_1\ne0$ gives $c=-1$.

The free propagator preserves tensor products.  Indeed, the Fourier transform
of $p(x_1)q(x_2)$ is $\widehat p(\xi_1)\widehat q(\xi_2)$, while
\[
 e^{-it(\xi_1^2+\xi_2^2)}=e^{-it\xi_1^2}e^{-it\xi_2^2}.
\]
Therefore
\[
 e^{it\Delta_2}(p\otimes q)
 =(e^{it\partial_x^2}p)\otimes(e^{it\partial_x^2}q).
\]

Define
\[
 A_t(x)=(1+2it)^{-1/2}
 \exp\!\left(-\frac{x^2}{2(1+2it)}\right),
\]
where the square root is chosen continuously from $t=0$.  Direct
differentiation gives
\[
 i\partial_tA_t+\partial_x^2A_t=0,
 \qquad A_0=\phi.
\]
Uniqueness of the free evolution yields $e^{it\partial_x^2}\phi=A_t$.
Since $\psi=-\phi'$ and spatial differentiation commutes with the free
propagator,
\[
 e^{it\partial_x^2}\psi
 =-\partial_xA_t=\frac{x}{1+2it}A_t.
\]
Consequently
\begin{align*}
 e^{it\Delta}f(x_1,x_2)
 &=A_t(x_1)A_t(x_2)\frac{x_2+ix_1}{1+2it},\\
 e^{it\Delta}g(x_1,x_2)
 &=A_t(x_1)A_t(x_2)\frac{x_2-ix_1}{1+2it}.
\end{align*}
Since $x_1,x_2$ are real,
\[
 |x_2+ix_1|^2=x_1^2+x_2^2=|x_2-ix_1|^2.
\]
All other factors in the two formulas are identical.  Their moduli therefore
coincide at every point and every time.

For $d>2$, tensor both functions with
$\Phi(x_3,\ldots,x_d)=\exp(-\frac12\sum_{j=3}^dx_j^2)$.  The free propagator
again factors, so the modulus identity and nonproportionality are preserved.
\end{proof}

\section{Comments on alternative proof methods}
\label{sec:stationary-comments}

The companion note contains three proofs for time-independent potentials.
They are recorded there because their hypotheses and methods differ from the
time-dependent argument of Theorem~\ref{thm:time-potential}.

\subsection{Masuda's unique-continuation theorem}

Assume that $V$ is time independent and that the two solutions belong to
$C(I;H^1(\R))$.  For
$\rho_w=|w|^2$ and
$j_w=\operatorname{Im}(\overline w\,\partial_xw)$, the equation with real
potential gives
\[
 \partial_t\rho_w+2\partial_x j_w=0.
\]
Equality of the densities therefore implies
$\partial_x(j_u-j_v)=0$.  The $H^1$ assumption gives
$\|j_w(t)\|_1\le \|w(t)\|_2\|\partial_xw(t)\|_2$, so the constant in $x$
must be zero.  Since
$2\operatorname{Re}(\overline w\,\partial_xw)=\partial_x|w|^2$, the real
parts agree as well, and hence
\[
 \overline u\,\partial_xu=\overline v\,\partial_xv
 \quad\text{almost everywhere on }I\times\R.
\]

The $H^1$ regularity also provides jointly continuous representatives.  On
the open set $\Omega=\{(t,x):u(t,x)\ne0\}$, define $q=v/u$.  Then $|q|=1$
and
\[
 0=\overline v\,\partial_xv-\overline u\,\partial_xu
   =|u|^2\overline q\,\partial_xq.
\]
Thus $q$ is constant in $x$ on each sufficiently small rectangle contained
in $\Omega$.  No division at a zero of $u$ is used.

The remaining point is to remove a possible dependence on time.  On such a
rectangle choose a nonzero, nonnegative $\chi\in C_c^\infty(\R)$ supported
in its spatial side, and write
\[
 A(t)=\int_\R\chi\overline u v\,dx,
 \qquad B(t)=\int_\R\chi|u|^2\,dx.
\]
If $v=c(t)u$ on the rectangle, then $A=cB$ and $B>0$.  Testing the two
equations and integrating once in $x$ gives $A'=cB'$ almost everywhere.
Consequently $c'B=0$, so $c$ is constant.  The quotient is therefore
constant on each connected component of $\Omega$.  For one such component,
$w=v-\zeta u$ is a solution which vanishes on a nonempty open subset of
spacetime.  Masuda's theorem~\cite{Masuda} yields $w\equiv0$.

This is the shortest of the three stationary proofs.  It was not used for the
main theorem because Masuda's result requires a time-independent semibounded
Hamiltonian, while the common potential arising in a nonlinear equation is
generally time dependent.  The phase reduction also uses $H^1$ regularity on
the observation interval.  Theorem~\ref{thm:time-potential}, by contrast,
allows the stated time-dependent potential class and arbitrary $L^2$ data.
The Masuda argument therefore has a narrower range for the applications of
this paper.

\subsection{Faddeev scattering}

For a real potential in $L^1_1(\R)$, one-dimensional scattering
theory~\cite{DeiftTrubowitz} gives finitely many simple negative eigenvalues
$-\kappa_j^2$ and two continuous channels.  The corresponding spectral
expansion has the form
\[
 e^{-itH}f
 =\sum_j c_j e^{it\kappa_j^2}\phi_j
  +\sum_{\alpha=1}^2\int_0^\infty
    e^{-itk^2}g_\alpha(k)e_\alpha(\,\cdot\,,k)\,dk.
\]
This representation remains valid in the resonant case at zero energy.  The
bound-state coefficients and the two functions $g_\alpha$ give a unitary
spectral representation of arbitrary $L^2$ initial data.

The temporal Fourier transform of the density identity separates products by
their energy difference $\delta=E-F$.  Along a fixed nonzero gap, use the
energy sum $s=E+F$.  If $y$ and $z$ solve the stationary equations at energies
$E$ and $F$, their products satisfy
\[
 A'=B,\qquad B'=(2V-s)A+2D,\qquad
 C'=-\delta A,\qquad
 D'=\left(V-\frac{s}{2}\right)B+\frac{\delta}{2}C,
\]
where $A=yz$, $B=y'z+yz'$, $C=y'z-yz'$, and $D=y'z'$.  The system contains
no derivative of $V$.  Moreover, when $E\ne F$, the four products obtained
from bases of the two two-dimensional solution spaces are linearly
independent: if their scalar component $A$ vanishes, the system successively
forces $B$, $D$, and $C$ to vanish.

To separate the energy sum, one divides a fixed-gap relation by $s-\lambda$
and integrates in $s$.  The resulting vector solves the same product system
with sum $\lambda$.  For sufficiently negative $\lambda$, both auxiliary
energies $(\lambda\pm\delta)/2$ are negative and the lower one lies below
$\inf\sigma(H)$.  Scattering asymptotics and the Riemann--Lebesgue lemma show
that its scalar component decays at both spatial infinities.  The
negative-energy ODE dichotomy then forces the whole vector to vanish: a
nonzero solution decaying at both ends at the lower energy would be an
$L^2$ eigenfunction below the spectrum.  Analyticity in $\lambda$ and
Stieltjes inversion show that the original fixed-gap measure is zero.
Applying this for almost every $\delta\ne0$, and treating the finite
bound--bound part separately, recovers
\[
 (\mathcal F u_0)\otimes\overline{\mathcal F u_0}
 =(\mathcal F v_0)\otimes\overline{\mathcal F v_0}.
\]
Equality of these rank-one operators gives $v_0=\zeta u_0$.

This proof allows arbitrary $L^2$ data, negative eigenvalues, and a
zero-energy resonance.  It was left in the companion note because it requires
a time-independent potential with Faddeev decay, observations for all
$t\in\R$, and the complete spectral resolution of $H$.  None of these
spectral reductions is available for the time-dependent potentials arising
in the nonlinear application.  The argument is therefore more specialized
than the physical-space proof of Theorem~\ref{thm:time-potential}.

\subsection{Observation on one exterior half-line}

Suppose that $V$ is smooth, compactly supported, and vanishes on
$(R,\infty)$.  Each negative-energy eigenfunction is then a nonzero multiple
of $e^{-\kappa_j(x-R)}$ on this half-line.  The two continuous scattering
channels can also be combined into one Fourier integral, so the restriction
of an evolution has the form
\[
 e^{-itH}f(x)
 =\sum_j a_j e^{-\kappa_j(x-R)}e^{it\kappa_j^2}
  +\int_\R F(k)e^{ik(x-R)-ik^2t}\,dk.
\]
Thus the half-line problem first reduces to phase retrieval for a free
Fourier extension coupled to finitely many decaying exponentials.

Test the density difference in time against the inverse Fourier transform of
$\chi\in C_c^\infty(\R)$, and take the Laplace transform in $y=x-R>0$.
If
\[
 q(k,\ell)=F(k)\overline{F(\ell)}
           -G(k)\overline{G(\ell)},
\]
then the free--free term becomes
\[
 \Phi_\chi(z)=\int_\R\frac{M_\chi(\xi)}{z-i\xi}\,d\xi,
 \qquad
 M_\chi(\xi)=\int_\R q(k,k-\xi)
                  \chi(\xi^2-2k\xi)\,dk,
 \qquad \operatorname{Re}z>0.
\]
The compact support of $\chi$ and a Schur estimate make $M_\chi$ integrable,
so this is an ordinary Cauchy transform.

Every term containing a bound state is holomorphic in a strip across the
imaginary axis.  The transformed density identity therefore continues
$\Phi_\chi$ through that axis.  The Cauchy jump formula gives
$M_\chi=0$ for every $\chi$.  The change of variables
$\xi=k-\ell$, $\omega=\ell^2-k^2$ then yields $q=0$ away from the diagonal,
which is a null set.  Hence $F\otimes\overline F=G\otimes\overline G$ and
$G=\zeta F$.  The mixed terms recover the bound-state coefficients by
successively crossing the lines $\operatorname{Re}z=-\kappa_j$; when the
continuous part vanishes, two elementary Vandermonde arguments recover the
finite bound-state rank-one matrix.  Finally, the transmission coefficient
has no positive-energy zeros, so the relation on the free half-line recovers
both original scattering channels.

This theorem uses less spatial data than Theorem~\ref{thm:time-potential},
but it assumes a time-independent, smooth, compactly supported potential and
observations for every time.  Its proof depends on an exactly free exterior
region, on analytic continuation of the associated Laplace transforms, and
on one-dimensional scattering coefficients.  These ingredients do not
persist for the time-dependent potentials in the nonlinear problem.  The
result was therefore retained as a separate auxiliary statement.

\begin{thebibliography}{99}

\bibitem{CarmeliHeinosaariSchultzToigo}
C.~Carmeli, T.~Heinosaari, J.~Schultz, and A.~Toigo,
\emph{Nonuniqueness of phase retrieval for three fractional Fourier
transforms},
Appl. Comput. Harmon. Anal. \textbf{39} (2015), 339--346.

\bibitem{DeiftTrubowitz}
P.~Deift and E.~Trubowitz,
\emph{Inverse scattering on the line},
Comm. Pure Appl. Math. \textbf{32} (1979), 121--251.

\bibitem{IonescuKenig}
A.~D.~Ionescu and C.~E.~Kenig,
\emph{$L^p$ Carleman inequalities and uniqueness of solutions of nonlinear
Schr\"odinger equations},
Acta Math. \textbf{193} (2004), 193--239.

\bibitem{Jaming2014}
P.~Jaming,
\emph{Uniqueness results in an extension of Pauli's phase retrieval problem},
Appl. Comput. Harmon. Anal. \textbf{37} (2014), 413--441.

\bibitem{Jaming2025}
P.~Jaming,
\emph{Phase retrieval for solutions of the Schr\"odinger equations},
in 2025 International Conference on Sampling Theory and Applications
(SampTA), 2025, 1--5.

\bibitem{JamingRathmair}
P.~Jaming and M.~Rathmair,
\emph{Uniqueness of phase retrieval from three measurements},
Adv. Comput. Math. \textbf{49} (2023), Paper No.~47.

\bibitem{KeelTao}
M.~Keel and T.~Tao,
\emph{Endpoint Strichartz estimates},
Amer. J. Math. \textbf{120} (1998), 955--980.

\bibitem{Masuda}
K.~Masuda,
\emph{A unique continuation theorem for solutions of the Schr\"odinger
equations},
Proc. Japan Acad. \textbf{43} (1967), 361--364.

\bibitem{MorozPerelomov}
B.~Z.~Moroz and A.~M.~Perelomov,
\emph{On a problem posed by Pauli},
Theoret. Math. Phys. \textbf{101} (1994), 1200--1204.

\bibitem{Pauli}
W.~Pauli,
\emph{Die allgemeinen Prinzipien der Wellenmechanik},
in H.~Geiger and K.~Scheel (eds.), Handbuch der Physik, Vol.~24,
Springer-Verlag, Berlin, 1933, 83--272.

\bibitem{TeschlBook}
G.~Teschl,
\emph{Mathematical Methods in Quantum Mechanics},
2nd ed., Graduate Studies in Mathematics 157,
American Mathematical Society, 2014.

\bibitem{VegaVisciglia}
L.~Vega and N.~Visciglia,
\emph{On the local smoothing for the Schr\"odinger equation},
Proc. Amer. Math. Soc. \textbf{135} (2007), 119--128.

\end{thebibliography}

\end{document}
